% Source PDF pages 38--40; manual pages 2-9--2-11.
\subsection{Geocentric Coordinate System}\label{sec:geocentric}

The geocentric coordinate system has its origin at the center of the earth, like the
ECFC system, but position and velocity are defined by magnitudes and orientation
angles. The position-vector angles, longitude $\lambda$ and geocentric latitude
$\delta_{gc}$, form a spherical coordinate system at the earth's center, as shown in
\cref{fig:geocentric-position-components}.

\begin{figure}[htbp]
  \centering
  \resizebox{0.62\linewidth}{!}{\begin{tikzpicture}[scale=0.98,>=Latex,line cap=round,line join=round]
  \coordinate (O) at (0,0);
  \coordinate (P) at (-0.95,2.95);
  \coordinate (Q) at (-0.95,-0.55);
  \draw[->,line width=0.9pt] (O) -- (-3.00,-1.55) node[below left] {$\hat{x}_{\oplus}$};
  \draw[->,line width=0.9pt] (O) -- (3.10,0) node[right] {$\hat{y}_{\oplus}$};
  \draw[->,line width=0.9pt] (O) -- (0,3.55) node[above] {$\hat{z}_{\oplus}$};
  \draw[->,line width=1.05pt] (O) -- (P) node[midway,right] {$\vec r$};
  \draw[dashed] (P) -- (Q);
  \draw[line width=0.9pt] (O) -- (Q);
  \fill (P) circle (1.6pt) node[left=5pt] {Vehicle};
  \fill (Q) circle (1.2pt);
  \node[align=left,right] at (1.10,-0.92) {Projection of $\vec r$ onto\\$\hat{x}_{\oplus}\hat{y}_{\oplus}$ equatorial plane};
  \draw[->] (0.50,-0.37) arc[start angle=-37,end angle=-143,radius=0.62];
  \node at (-0.10,-0.70) {$\lambda$};
  \draw[->] (-0.48,0.14) arc[start angle=165,end angle=104,radius=0.55];
  \node at (-0.80,0.48) {$\delta_{gc}$};
\end{tikzpicture}
}
  \caption{Geocentric position components.}
  \label{fig:geocentric-position-components}
\end{figure}

\taossubhead{Geocentric Position}

\noindent\begin{tabularx}{\textwidth}{@{}>{\ttfamily}l>{$\displaystyle}l<{$}X@{}}
  rcm & \lVert\vec{r}\rVert &
    Magnitude of the position vector, or distance from the earth's center. \\
  long & \lambda &
    Longitude: the angle between the ECFC $x$ axis and the projection of
    $\vec{r}$ onto the equatorial plane. It is positive eastward, measured
    counterclockwise from the ECFC $x$ axis. \\
  latgc & \delta_{gc} &
    Geocentric latitude: the angle between $\vec{r}$ and the equatorial plane,
    positive northward. \\
\end{tabularx}

The function \texttt{p\_geocentric\_to\_ecfc} uses
\begin{equation}
\begin{aligned}
  \vec{r}\cdot\hat{x}_{\oplus}
    &= \lVert\vec{r}\rVert\cos\delta_{gc}\cos\lambda,\\
  \vec{r}\cdot\hat{y}_{\oplus}
    &= \lVert\vec{r}\rVert\cos\delta_{gc}\sin\lambda,\\
  \vec{r}\cdot\hat{z}_{\oplus}
    &= \lVert\vec{r}\rVert\sin\delta_{gc}.
\end{aligned}
\label{eq:geocentric-position-to-ecfc}
\end{equation}

The reverse relationships are evaluated by
\texttt{p\_ecfc\_to\_geocentric}:
\begin{equation}
  \tan\lambda
    = \frac{\vec{r}\cdot\hat{y}_{\oplus}}
           {\vec{r}\cdot\hat{x}_{\oplus}},
  \label{eq:longitude-from-ecfc}
\end{equation}
\begin{equation}
  \sin\delta_{gc}
    = \frac{\vec{r}\cdot\hat{z}_{\oplus}}
           {\lVert\vec{r}\rVert}.
  \label{eq:geocentric-latitude-from-ecfc}
\end{equation}

\taossubhead{Geocentric Velocity}

The velocity orientation angles are measured relative to the local geocentric horizon
system described in \cref{sec:local-geocentric}. As shown in
\cref{fig:geocentric-velocity-components}, they are analogous to longitude and
latitude, but are referenced to the local horizon rather than to ECFC.

\begin{figure}[htbp]
  \centering
  \resizebox{0.64\linewidth}{!}{\begin{tikzpicture}[scale=1.02,>=Latex,line cap=round,line join=round]
  \coordinate (O) at (0,0);
  \coordinate (V) at (1.15,2.65);
  \coordinate (H) at (1.75,0.62);
  \draw[->,line width=0.9pt] (O) -- (-1.75,1.90) node[above left] {$\hat{x}_{gc}$};
  \draw[->,line width=0.9pt] (O) -- (3.00,0) node[right] {$\hat{y}_{gc}$};
  \draw[->,line width=0.9pt] (O) -- (0,-1.60) node[below] {$\hat{z}_{gc}$};
  \draw[->,line width=1.05pt] (O) -- (V) node[above left] {$\vec V_{\oplus}$};
  \draw[line width=0.85pt] (O) -- (H);
  \draw[dashed] (V) -- (H);
  \node[align=left,right] at (2.12,1.40) {Projection of $\vec V_{\oplus}$\\onto $\hat{x}_{gc}\hat{y}_{gc}$ plane};
  \draw[->] (0.75,0.27) arc[start angle=20,end angle=67,radius=0.82];
  \node at (1.08,0.95) {$\gamma_{gc}$};
  \draw[->] (-0.32,0.48) arc[start angle=124,end angle=28,radius=0.58];
  \node at (0.10,0.73) {$\psi_{gc}$};
  \node[below left] at (O) {Vehicle};
\end{tikzpicture}
}
  \caption{Geocentric velocity components.}
  \label{fig:geocentric-velocity-components}
\end{figure}

\noindent\begin{tabularx}{\textwidth}{@{}>{\ttfamily}l>{$\displaystyle}l<{$}X@{}}
  vel & \lVert\vec{V}_{\oplus}\rVert &
    Magnitude of the earth-relative velocity vector. \\
  gamgc & \gamma_{gc} &
    Geocentric vertical flight-path angle: the angle between
    $\vec{V}_{\oplus}$ and the local geocentric horizon plane. It is positive
    when the vehicle moves away from the earth's surface. \\
  psigc & \psi_{gc} &
    Geocentric horizontal flight-path, or heading, angle: the angle between local
    north $\hat{x}_{gc}$ and the projection of $\vec{V}_{\oplus}$ onto the local
    horizon plane. It is positive clockwise from north. \\
\end{tabularx}

The velocity components in local geocentric coordinates are
\begin{equation}
  \vec{V}_{\oplus}\cdot\hat{x}_{gc}
    = \lVert\vec{V}_{\oplus}\rVert
      \cos\gamma_{gc}\cos\psi_{gc},
  \label{eq:geocentric-velocity-x}
\end{equation}
\begin{equation}
  \vec{V}_{\oplus}\cdot\hat{y}_{gc}
    = \lVert\vec{V}_{\oplus}\rVert
      \cos\gamma_{gc}\sin\psi_{gc},
  \label{eq:geocentric-velocity-y}
\end{equation}
\begin{equation}
  \vec{V}_{\oplus}\cdot\hat{z}_{gc}
    = -\lVert\vec{V}_{\oplus}\rVert\sin\gamma_{gc}.
  \label{eq:geocentric-velocity-z}
\end{equation}
Solving for the flight-path angles gives
\begin{equation}
  \sin\gamma_{gc}
    = -\frac{\vec{V}_{\oplus}\cdot\hat{z}_{gc}}
             {\lVert\vec{V}_{\oplus}\rVert},
  \label{eq:geocentric-gamma-from-ecfc}
\end{equation}
\begin{equation}
  \tan\psi_{gc}
    = \frac{\vec{V}_{\oplus}\cdot\hat{y}_{gc}}
           {\vec{V}_{\oplus}\cdot\hat{x}_{gc}}.
  \label{eq:geocentric-psi-from-ecfc}
\end{equation}
These quantities are calculated in \texttt{v\_ecfc\_to\_geocentric}. If
$\vec{V}_{\oplus}=\vec{0}$, both flight-path angles are undefined and TAOS sets
$\gamma_{gc}=\psi_{gc}=0$. If the velocity is straight up or down
($\lvert\gamma_{gc}\rvert=90^{\circ}$), heading is undefined and TAOS sets it to
zero.

The reverse calculation, performed by
\texttt{v\_geocentric\_to\_ecfc}, is
\begin{equation}
\begin{aligned}
  \vec{V}_{\oplus}\cdot\hat{x}_{\oplus}
    ={}& -(\vec{V}_{\oplus}\cdot\hat{x}_{gc})
          \cos\lambda\sin\delta_{gc}
        -(\vec{V}_{\oplus}\cdot\hat{y}_{gc})\sin\lambda\\
       &-(\vec{V}_{\oplus}\cdot\hat{z}_{gc})
          \cos\lambda\cos\delta_{gc},
\end{aligned}
\label{eq:geocentric-to-ecfc-x}
\end{equation}
\begin{equation}
\begin{aligned}
  \vec{V}_{\oplus}\cdot\hat{y}_{\oplus}
    ={}& -(\vec{V}_{\oplus}\cdot\hat{x}_{gc})
          \sin\lambda\sin\delta_{gc}
        +(\vec{V}_{\oplus}\cdot\hat{y}_{gc})\cos\lambda\\
       &-(\vec{V}_{\oplus}\cdot\hat{z}_{gc})
          \sin\lambda\cos\delta_{gc},
\end{aligned}
\label{eq:geocentric-to-ecfc-y}
\end{equation}
\begin{equation}
  \vec{V}_{\oplus}\cdot\hat{z}_{\oplus}
    = (\vec{V}_{\oplus}\cdot\hat{x}_{gc})\cos\delta_{gc}
      -(\vec{V}_{\oplus}\cdot\hat{z}_{gc})\sin\delta_{gc}.
  \label{eq:geocentric-to-ecfc-z}
\end{equation}
The local components in these equations are supplied by
\cref{eq:geocentric-velocity-x,eq:geocentric-velocity-y,eq:geocentric-velocity-z}.
